\setcounter{secnumdepth}{2}


\appendix 


\section{Implementation Details and Efficiency Analysis}
\label{sec:Evo-Merging Implement Detail}

\subsection{Model Preparation and Inference Environment}
To construct our diverse source pool of over 100 specialist models, we fine-tuned each candidate using the AdamW optimizer with a batch size of 12, an initial learning rate of 1e-4, and a cosine scheduler featuring a 0.1 warm-up ratio. 
During the evaluation and API-based feedback collection, we employed a nucleus sampling strategy with a temperature of 0.95, a top\_p of 0.7, and a top\_k of 50 to ensure generation diversity.
All inference processes were deployed on an NVIDIA A800 GPU (80 GB VRAM), while the core evolutionary optimization (Evo-Merging) was executed on a standard Intel(R) Xeon(R) Platinum 8352Y CPU @ 2.20GHz.

\subsection{Evolutionary Search Configuration}
We tailored the search configuration to the complexity of the task and the size of the model pool. 
To mitigate overfitting during the derivative-free optimization, we employed a dynamic early-stopping strategy based on a minimal validation set. The specific convergence patterns observed are as follows: for pools with fewer than 10 models, convergence in Stage 2 is typically achieved within 10 generations; for 25--50 models, it requires 10--20 generations; and for massive pools ($>$50 models), convergence occurs within 20--40 generations. The detailed setup for our two primary experimental settings is summarized in Table~\ref{tab:exp_setups}.

\begin{table}[htbp]
\centering
\small
\caption{Experimental configurations for Evolutionary Optimization.}
\label{tab:exp_setups}
\vspace{-2mm}
\setlength{\tabcolsep}{2.5mm}
\begin{tabular}{lcc}
\toprule
\textbf{Configuration} & \textbf{OOD Setting} & \textbf{ID Setting} \\ 
\midrule
Model Pool Size ($N$) & 100+ & $N < 10$ \\ 
Validation Set Size & 300 samples (Total) & 100 samples/task \\ 
Stage 1 Max Iterations & 20 & 20 \\ 
Stage 2 Max Iterations & 40 & 20 \\ 
\bottomrule
\end{tabular}
\vspace{-2mm}
\end{table}

\subsection{Computational Efficiency Analysis}
A practical advantage of \texttt{Evo-Merging} is that it separates the one-time offline merge search from online serving. The method is derivative-free and uses restricted black-box access, but we do not claim a formal privacy guarantee.

\paragraph{Optimized Parameter Count.} 
Unlike gradient-based methods (e.g., LoRAHub) that may require back-propagating through high-dimensional weights, our framework operates on a significantly reduced parameter space. For a pool of $N$ models, we optimize only two scalar vectors: the sparsity vector $\bm{\alpha} \in \mathbb{R}^N$ and the scaling vector $\bm{\beta} \in \mathbb{R}^N$. Thus, the total number of optimized parameters is merely $2N$. Even for our massive repository ($N=100$), this results in only 200 learnable scalars, ensuring rapid convergence and stability.

\paragraph{Resource Consumption and Training Time.} 
Since the optimization process is derivative-free, it eliminates the need to store gradient graphs or optimizer states (e.g., momentum) in VRAM. The evolutionary search is computationally lightweight and runs entirely on the CPU. The dominant offline cost is repeated validation-time inference rather than the optimizer itself. We decompose the merge-time cost as
\begin{equation}
T_{\mathrm{merge}} = T_{\mathrm{opt}} + n_{\mathrm{eval}} \cdot T_{\mathrm{val}},
\end{equation}
where $T_{\mathrm{opt}}$ is the CPU-side optimizer overhead, $T_{\mathrm{val}}$ is the time required for one validation evaluation, and $n_{\mathrm{eval}}$ is determined by population size, validation-set size, and the number of generations. For the OOD setting, the population size is 20, the validation set contains 300 samples, and the maximum Stage 1/Stage 2 iterations are 20/40. Under the rebuttal-time accounting used for this setting, this corresponds to roughly $20 \times 300 \times 40 = 240{,}000$ batchable validation inferences. After this offline search, deployment uses a single merged adapter for inference.


\section{Qwen detail}
\label{sec:Qwen detail}
To further validate our approach, we adopted the powerful Qwen model and evaluated on the InstructUIE benchmark~\cite{wang2023instructuiemultitaskinstructiontuning}, where Qwen outperforms alternatives like Llama3. This stronger base model improved all methods, including baselines.
Nevertheless, \texttt{Evo-Merging} remains consistently superior. As shown in Table~\ref{tab:qwen_performance_indomain}, it achieves the highest average Precision and F1 on in-domain tasks. Similarly, on out-of-domain (OOD) tasks (Table~\ref{tab:method_performance_ood_qwen}), it again attains the top average F1-score, confirming its effectiveness even with a more advanced base model.




\begin{table*}[t!]
    \caption{Full per-task ID results for Qwen2.5 7B, corresponding to the ID averages in Table~\ref{tab:qwen_merged_results}. $\dagger$ indicates that the improvement of \texttt{Evo-Merging} over the strongest baseline (LoRAHub) is statistically significant ($p < 0.05$).}
  \label{tab:qwen_performance_indomain}
  \vspace{-2mm}
  \centering
  \setlength{\tabcolsep}{1.3pt}
  \begin{tabular}{l *{9}{cc}}
    \toprule
    \multirow{2}{*}{\textbf{Method}}
    & \multicolumn{2}{c}{ACE\_2005} & \multicolumn{2}{c}{CrossNER\_AI} & \multicolumn{2}{c}{CrossNER\_Lit} & \multicolumn{2}{c}{FabNER} & \multicolumn{2}{c}{Mit-Movie} & \multicolumn{2}{c}{MultiNERD} & \multicolumn{2}{c}{TweetNER} & \multicolumn{2}{c}{WikiANN} & \multicolumn{2}{c}{Avg} \\
    \cmidrule(lr){2-3} \cmidrule(lr){4-5} \cmidrule(lr){6-7} \cmidrule(lr){8-9} \cmidrule(lr){10-11} \cmidrule(lr){12-13} \cmidrule(lr){14-15} \cmidrule(lr){16-17} \cmidrule(lr){18-19}
    & Prec & F1 & Prec & F1 & Prec & F1 & Prec & F1 & Prec & F1 & Prec & F1 & Prec & F1 & Prec & F1 & Prec & F1 \\
    \midrule
    Individual model & 74.74 & 74.05 & 37.31 & 34.11 & 44.54 & 41.96 & 52.86 & 50.25 & 83.65 & 82.86 & 89.06 & 88.70 & 60.68 & 58.00 & 84.89 & 84.67 & 65.97 & 64.32 \\
    \hline
    TIES & 27.32 & 21.71 & 43.25 & 39.67 & \textbf{43.85} & \textbf{40.41} & 25.20 & 19.79 & 58.27 & 56.62 & 53.00 & 60.14 & 35.61 & 36.11 & 56.38 & 56.28 & 42.86 & 41.34 \\
    DARE & 26.98 & 21.53 & 43.82 & 40.23 & 43.74 & 40.26 & 26.10 & 20.54 & 58.52 & 56.89 & 53.29 & 60.41 & 35.27 & 35.98 & 56.70 & 56.69 & 43.05 & 41.57 \\
    DELLA & 26.93 & 21.43 & 43.77 & 40.22 & 43.55 & 39.99 & 26.37 & 20.75 & 58.98 & \textbf{57.32} & 53.20 & 60.25 & 35.31 & 36.05 & 56.67 & 56.58 & 43.10 & 41.57 \\
    Breadcrumbs & 26.40 & 21.11 & 42.09 & 38.87 & 42.58 & 39.58 & 23.98 & 18.92 & 57.21 & 55.78 & 51.75 & 58.61 & 34.78 & 35.30 & 54.89 & 54.85 & 41.71 & 40.38 \\
    Task Arithmetic & 26.83 & 21.29 & 43.17 & 39.63 & 43.54 & 40.11 & 25.48 & 20.01 & 57.84 & 56.27 & 53.15 & 60.24 & 34.69 & 35.12 & 55.75 & 55.79 & 42.56 & 41.06 \\
    LoRAHub & \textbf{31.17} & \textbf{23.24} & 35.79 & 28.77 & 35.81 & 30.32 & 17.31 & 12.11 & 54.94 & 49.17 & 60.46 & 64.58 & 41.82 & 36.92 & 64.09 & 62.26 & 42.67 & 38.42 \\
    EMR-Merging & 27.19 & 21.78 & 43.66 & 40.38 & 42.76 & 39.52 & 25.55 & 20.05 & 57.77 & 56.33 & 52.71 & 59.91 & 34.71 & 35.52 & 55.66 & 55.71 & 42.50 & 41.15 \\
    \midrule
    \textbf{\texttt{Evo-Merging}} & 25.90 & 19.62 & \textbf{49.02} & \textbf{44.08} & 39.14 & 34.84 & \textbf{28.08} & \textbf{21.34} & \textbf{62.67} & 55.46 & \textbf{62.55} & \textbf{67.41} & \textbf{43.52} & \textbf{41.25} & \textbf{64.59} & \textbf{63.10} & \textbf{46.93} & \textbf{43.39}$^\dagger$ \\
    \bottomrule
  \end{tabular}
  \vspace{-2mm}
\end{table*}



\begin{table*}[t!]
  \caption{Full per-task OOD results for Qwen2.5 7B, corresponding to the OOD averages in Table~\ref{tab:qwen_merged_results}.}
  \label{tab:method_performance_ood_qwen}
  \vspace{-2mm}
  \setlength{\tabcolsep}{4.0pt}
  \begin{tabular}{l cc cc cc cc cc cc c c}
    \toprule
    \multirow{3}{*}{Method}  & \multicolumn{4}{c}{NER} & \multicolumn{4}{c}{RE} & \multicolumn{4}{c}{ET} & \multicolumn{2}{c}{\multirow{2}{*}{Avg}} \\ 
    \cmidrule(lr){2-5} \cmidrule(lr){6-9} \cmidrule(lr){10-13}
     & \multicolumn{2}{c}{Mit-Restaurant} & \multicolumn{2}{c}{FindVehicle} & \multicolumn{2}{c}{New York Time} & \multicolumn{2}{c}{Conll04} & \multicolumn{2}{c}{Mit-Movie} & \multicolumn{2}{c}{FabNER} & & \\
    \cmidrule(lr){2-3} \cmidrule(lr){4-5} \cmidrule(lr){6-7} \cmidrule(lr){8-9} \cmidrule(lr){10-11} \cmidrule(lr){12-13} \cmidrule(lr){14-15}
      & Prec & F1 & Prec & F1 & Prec & F1 & Prec & F1 & Prec & F1 & Prec & F1 & Prec & F1 \\
    \midrule
    TIES            & 61.78 & 48.34 & 23.81 & 20.98 & 32.64 & 29.79 & 19.70 & 17.77 & 58.47 & 47.13 & 50.91 & 33.17 & 41.22 & 32.86 \\
    DARE (TIES+)    & 58.97 & 45.75 & 16.65 & 14.29 & 31.95 & 29.28 & 18.83 & 16.96 & 50.51 & 39.88 & 48.98 & 31.72 & 37.65 & 29.65 \\
    Task Arithmetic & 63.52 & 50.65 & 29.49 & 25.99 & 25.28 & 23.24 & 19.70 & 17.82 & 59.60 & 47.08 & 50.50 & 32.21 & 41.35 & 32.83 \\
    Breadcrumbs     & 63.15 & 50.89 & 42.94 & 39.68 & 24.21 & 22.01 & 22.29 & 19.91 & 62.61 & 51.18 & 49.58 & 32.43 & 44.13 & 36.02 \\
    Della           & 59.05 & 45.82 & 16.93 & 14.58 & 31.30 & 28.45 & 18.40 & 16.52 & 51.56 & 40.81 & 49.05 & 31.78 & 37.72 & 29.66 \\
    EMR-Merging     & 60.82 & 48.18 & 24.78 & 21.32 & 43.16 & 38.60 & 19.70 & 17.77 & 58.23 & 45.51 & 50.41 & 31.15 & 42.85 & 33.76 \\

    Best Adapter    & 63.66 & 57.64 & 59.48 & 56.67 & 63.23 & 55.52 & 33.12 & 31.21 & 64.15 & 63.72 & \textbf{53.07} & \textbf{50.28} & 56.12 & 52.51 \\
    LoRA            & 59.27 & 56.83 & 52.87 & 51.62 & 24.32 & 22.75 & 33.12 & 31.21 & 59.27 & 56.83 & 11.13 & 10.70 & 40.00 & 38.32 \\
    LoRAHub         & 56.20 & 52.13 & \textbf{71.64} & \textbf{70.86} & \textbf{61.17} & \textbf{61.15} & 31.23 & 29.77 & \textbf{77.96} & \textbf{77.06} & 45.45 & 46.72 & 57.28 & 56.28 \\
    \midrule
    \textbf{\texttt{Evo-Merging}} & \textbf{70.86} & \textbf{69.55} & 69.70 & 69.36 & 60.24 & 58.95 & \textbf{47.29} & \textbf{44.19} & 75.20 & 75.36 & 51.39 & 49.97 & \textbf{62.45} & \textbf{61.23} \\
    \bottomrule
  \end{tabular}
  \vspace{-2mm}
\end{table*}


\section{Dataset Details}
\label{sec: Dataset Details and Evalution}

Our methodology is grounded in the \textbf{InstructUIE} benchmark, which unifies diverse Information Extraction (IE) tasks into a single text-to-text format. We created a heterogeneous model pool by fine-tuning over 100 specialist models, each mastering one of the nine distinct IE task types defined by the framework. These tasks are structured in a three-layer hierarchy.

\subsection*{Hierarchy of Information Extraction Tasks}

Our source pool contains specialists for all tasks across three hierarchical layers of increasing complexity:
1) Layer 1: Named Entity Recognition. The primary \textbf{NER} task (identifying an entity's span and type) is decomposed into its auxiliary components: Entity Span Extraction (ES) for boundaries and Entity Typing (ET) for classification.
2) Layer 2: Relation Extraction. The primary \textbf{RE} task (extracting a semantic triplet) is decomposed into identifying a related Entity Pair (EP) and classifying the Entity Pair Relationship (EPR).
3) Layer 3: Event Extraction. The complex primary \textbf{EE} task (extracting an event with its trigger and arguments) is broken down into Event Trigger Identification (EET) and Event Argument Extraction (EEA).

This framework is designed for two primary use cases: it can either fuse the entire model pool to create a model that generalizes effectively to new, out-of-domain (OOD) tasks, or selectively combine a subset of expert models to create a proficient multi-task specialist for their specific in-domain (ID) tasks.

\subsection*{Model Category}
To analyze how different source models contribute to a target task, we classify them into five categories. We use the \textbf{NER FindVehicle} task as a case study, as illustrated in Figure~\ref{fig:fusion_solution_analysis}. The categories are as follows:
1) \textbf{Target Domain:} Models trained on the target dataset itself. For the FindVehicle task, these are the \texttt{ES\_FindVehicle} and \texttt{ET\_FindVehicle} models. As auxiliary tasks, they are highly specialized and, as shown in Figure~\ref{fig:fusion_solution_analysis}, provide a significant positive contribution.
2) \textbf{Synergistic Task:} Models trained on the same task type (e.g., NER, ET, ES) but on different, general-domain datasets. These models learn universal features of the task that are beneficial for the target.
3) \textbf{Irrelevant Domain (Bio/Chem):} Specialist models trained on domains with distinct vocabularies and contexts, such as biology and chemistry. These are expected to contribute negatively or be pruned heavily.
4) \textbf{Irrelevant Task (RE/EE):} Models trained on fundamentally different tasks, such as Relation Extraction (RE) and Event Extraction (EE), which belong to a different branch of the task hierarchy.
5) \textbf{Other:} A broader category for models trained on auxiliary tasks (ET, ES) from domains that are neither clearly synergistic nor explicitly irrelevant.

\section{Further Analysis of Black-Box Merging}
\label{sec:more_analysis}

To further contextualize the advantages of \texttt{Evo-Merging}, we analyze its design principles against three distinct categories of methods. This analysis highlights why our two-stage framework is uniquely suited for the challenging setting of Black-Box Merging (BMM) with massive model repositories.

\paragraph{Comparison with White-box Task Vector Methods.}
Methods such as TIES-Merging, DARE, and Task Arithmetic typically assume direct parameter access (White-box). While they demonstrate considerable success in resolving conflicts within controlled In-Domain settings, they exhibit a critical structural limitation in the BMM context: \textbf{susceptibility to signal dilution}. 
They lack an intrinsic mechanism to filter out irrelevant or detrimental models from a massive, noisy pool. Consequently, their application in such settings often necessitates laborious manual pre-selection of a small, promising subset of models and is highly sensitive to hyperparameter tuning. This theoretical limitation explains the \textbf{catastrophic performance decline} observed when they confront a large-scale repository (e.g., 100+ models). It is precisely this vulnerability that motivates the pre-selection strategy detailed in Appendix~\ref{sec: Pre-selection Strategy for Baselines in OOD Setting}, whereas our automated two-stage framework is designed to work without such assistance.

\paragraph{Comparison with Black-box Compatible Composition.}
Dynamic methods like LoRA-Flow~\cite{wang2024flow} are incompatible with LMaaS, as they require gradient feedback, violating the strict black-box constraint. 
Although the gradient-free \textbf{LoRAHub} serves as the primary Black-box compatible baseline, its performance degrades significantly with large, noisy model pools. This is largely due to the absence of a dedicated denoising mechanism to handle the exponentially complex search space of 100+ models. 
In contrast, our framework's \textbf{Stage 1 (Sparsity-Based Denoising)} directly addresses this limitation by proactively identifying and pruning noisy parameters before fusion, ensuring robust performance even at a massive scale.

\paragraph{Comparison with Single-Source PEFT Methods.}
Standard PEFT methods like LoRA and (IA)\textsuperscript{3} efficiently adapt a single model to new tasks with minimal updates. However, they underperform in few-shot OOD generalization compared to merging approaches, as they are limited by their reliance on a single pre-trained source. 
In contrast, \texttt{Evo-Merging} leverages the \textbf{collective knowledge} across a diverse repository of expert models. By integrating the rich "parametric memory" distributed across many models, our framework achieves superior few-shot generalization and emergent capabilities—advantages that are theoretically beyond the scope of adapting a single model instance.



\section{Pre-selection Strategy for Baselines}
\label{sec: Pre-selection Strategy for Baselines in OOD Setting}
\paragraph{Motivation for Pre-selection.} 
The large (over 100), noisy model pool in the out-of-domain (OOD) setting poses a significant challenge for most baseline methods, leading to performance degradation and making hyperparameter tuning infeasible. We use the following two-stage pre-selection strategy for these baselines: 1) Model Pre-selection: For each target OOD task, we first identified a promising subset of source models (typically less than 10) based on their generalization performance on the same validation set. This step filters for the most relevant models. The resulting curated adapter sets used for each task are detailed in Table~\ref{tab:model_distribution_standardized}.
2) Hyperparameter Tuning: On these manageable subsets, we then performed an extensive grid search to find the optimal hyperparameters for each baseline (e.g., density and $\lambda$). The best-performing configurations, which were used in our main results, are summarized in Table~\ref{sec:appendix_hyperparams}.
This pre-selection strategy was not applied to \texttt{Evo-Merging} or LoRAHub; both are evaluated on the full repository of 100+ models. We additionally ran \texttt{Evo-Merging} on the same oracle top-$K$ subsets used by this strategy. As shown in Table~\ref{tab:evo_oracle_topk}, \texttt{Evo-Merging} reaches 53.68\% average F1, slightly above its full-repository result of 52.13\%.




\section{More Analysis About Evaluation of Cross-task/domain Robustness }
\label{sec: More analyse about Evo-Merging on Muti-task merging}
To rigorously evaluate the robustness of our proposed merging algorithm, we introduce a set of five distractor tasks into the experimental setup. These tasks are intentionally selected from the Cross-Task benchmark, namely Relation Extraction (\texttt{RE\_NYT11}) and Event Extraction (\texttt{EE\_PHEE}, \texttt{EE\_CASIE}), to serve as sources of interference. Additionally, we incorporate two Entity Span Extraction tasks (\texttt{ES\_ACE\_2004}, \texttt{ES\_ACE\_2005}) to act as auxiliary tasks for Named Entity Recognition (NER). This setup is designed to stress-test the method's ability to handle parameter competition and mitigate negative transfer from cross-task learning.

The results presented in Table~\ref{tab:method_performance_indomain_Robustness} highlight a critical challenge in in-domain merging. Most baseline methods exhibit a significant performance degradation in the presence of the distractor tasks, a phenomenon akin to catastrophic forgetting. This decline can be largely attributed to the interference from heterogeneous task types, conflicting entity label sets, and the semantic divergence introduced by unrelated domains. In stark contrast, our Evo-Merging approach demonstrates remarkable resilience. By leveraging our specialized denoising and sign-flip scaling mechanisms, Evo-Merging effectively mitigates the negative influence from these unrelated tasks. It successfully filters parameter-level noise and harmonizes conflicting parameter directions, allowing it to distill shared, underlying knowledge from the diverse model pool. Consequently, rather than suffering from interference, our method thrives in scenarios with a greater number of models, achieving superior performance by leveraging the increased diversity.



\clearpage


\begin{table*}[htbp]
\centering
\caption{\texttt{Evo-Merging} evaluated on the same oracle top-$K$ OOD subsets used by white-box baselines; this diagnostic table supports the pre-selection discussion in Appendix~\ref{sec: Pre-selection Strategy for Baselines in OOD Setting}.}
\label{tab:evo_oracle_topk}
\vspace{-2mm}
\setlength{\tabcolsep}{4pt}
\begin{tabular}{lccccccc}
\toprule
Method & ET\_mit-movie & NER\_FindVehicle & ET\_FabNER & NER\_mit-restaurant & RE\_NYT & RE\_conll04 & Avg F1 \\
\midrule
\texttt{Evo-Merging} (oracle top-$K$) & 88.00 & 59.33 & 37.76 & 46.68 & 54.52 & 35.77 & 53.68 \\
\bottomrule
\end{tabular}
\vspace{-2mm}
\end{table*}



\begin{table*}[htbp]
\centering
\caption{Summary of hyperparameter grid search for baseline methods used in the OOD pre-selection protocol.}
\label{tab:grid_search_summary}
\vspace{-2mm}
\setlength{\tabcolsep}{2.5pt}
\begin{tabular}{l|ccc|ccc|ccc|ccc}
\toprule
\multirow{2}{*}{\textbf{Dataset}} & \multicolumn{3}{c|}{\textbf{TIES-Merging}} & \multicolumn{3}{c|}{\textbf{DARE}} & \multicolumn{3}{c|}{\textbf{Della}} & \multicolumn{3}{c}{\textbf{Breadcrumbs}} \\
\cmidrule(lr){2-4} \cmidrule(lr){5-7} \cmidrule(lr){8-10} \cmidrule(l){11-13}
& \textbf{Params ($\lambda, p$)} & \textbf{Prec.} & \textbf{F1} & \textbf{Params ($\lambda, p$)} & \textbf{Prec.} & \textbf{F1} & \textbf{Params ($\lambda, p, \epsilon$)} & \textbf{Prec.} & \textbf{F1} & \textbf{Params ($\lambda, p, \gamma$)} & \textbf{Prec.} & \textbf{F1} \\
\midrule
NER-mitrestaurant & (1.0, 0.9) & 21.40 & 24.22 & (1.0, 0.9) & 10.88 & 14.31 & (0.5, 0.9, 0.1) & 29.69 & 26.06 & (0.5, 0.7, 0.01) & 38.99 & 34.06 \\
NER-FindVehicle   & (1.0, 0.9) & 29.87 & 29.45 & (0.3, 0.9) & 24.13 & 23.69 & (0.7, 0.9, 0.1) & 16.66 & 17.04 & (0.5, 0.9, 0.01) & 47.41 & 38.42 \\
ET-mitmovie       & (0.9, 0.9) & 48.67 & 51.12 & (0.1, 0.5) & 44.57 & 48.60 & (0.7, 0.9, 0.1) & 43.85 & 47.68 & (0.5, 0.7, 0.01) & 49.73 & 51.60 \\
ET-FabNER         & (0.5, 0.9) & 15.91 & 11.42 & (0.1, 0.9) & 15.44 & 10.71 & (1.0, 0.9, 0.1) & 14.41 & 12.51 & (0.5, 0.9, 0.01) & 20.53 & 16.84 \\
RE-NYT            & (1.0, 0.9) & 12.55 & 7.57  & (0.1, 0.9) & 3.69  & 2.53  & (0.5, 0.9, 0.1) & 3.94  & 3.25  & (0.5, 0.9, 0.01) & 34.05 & 28.03 \\
RE-CoNLL04        & (1.0, 0.9) & 20.78 & 19.05 & (0.1, 0.9) & 10.85 & 8.08  & (1.0, 0.9, 0.1) & 19.91 & 17.97 & (0.5, 0.9, 0.01) & 26.30 & 24.56 \\
\bottomrule
\end{tabular}
\vspace{-2mm}
\end{table*}

\begin{table*}[htbp]
\small
\centering
\caption{Standardized model distribution across datasets. This table lists the top-performing source adapters selected by the pre-selection process for each target OOD task.}
\label{tab:model_distribution_standardized}
\vspace{-2mm}
\small
\setlength{\tabcolsep}{1pt}
\begin{tabular}{llllll} 
\toprule
\textbf{Target: MIT-Restaurant} & \textbf{Target: FindVehicle} & \textbf{Target: New York Times} & \textbf{Target: CoNLL04} & \textbf{Target: MIT-Movie} & \textbf{Target: FabNER} \\
\midrule
NER CrossNER AI         & NER MIT-Movie SFT       & NER CoNLL 2003     & EPR CoNLL04         & NER MIT-Movie           & ET MultiNERD \\
NER CrossNER Literature & NER CoNLL 2003          & EPR New York Times & EPR SciERC          & NER MIT-Restaurant      & ET CoNLL 2003 \\
NER CrossNER Music      & NER OntoNotes           & RE NYT11           & EPR New York Times  & NER FindVehicle         & ET MIT-Movie \\
NER CrossNER Politics   & NER CrossNER Politics   &                    & RE SciERC           & NER FabNER              & ET CrossNER Science \\
NER CrossNER Science    & NER Broad Tweet Corpus  &                    &                     & NER CrossNER Literature & ET Broad Tweet Corpus \\
NER FabNER              & NER WikiANN EN          &                    &                     & ET CrossNER AI          & ET OntoNotes \\
NER FindVehicle         & NER CrossNER Science    &                    &                     & NER CrossNER AI         & ET WikiANN EN \\
NER MIT-Movie           & NER TweetNER7           &                    &                     & ET CrossNER Music       & ET WikiNeural \\
NER TweetNER7           & NER ACE 2005            &                    &                     & NER CrossNER Music      & ET FindVehicle \\
NER WikiANN EN          & NER MultiNERD           &                    &                     & ET CrossNER Literature  & ET MIT-Restaurant \\
\bottomrule
\end{tabular}
\vspace{-2mm}
\end{table*}

\begin{table*}[htbp]
\centering
\caption{Main Results for the Evaluation of In-Domain (ID) Robustness.}
\label{tab:method_performance_indomain_Robustness}
\vspace{-2mm}
\setlength{\tabcolsep}{1.8pt}
\begin{tabular}{l *{9}{cc}}
\toprule
\multirow{2}{*}{\textbf{Method}}
& \multicolumn{2}{c}{ACE\_2005} & \multicolumn{2}{c}{CrossNER\_AI} & \multicolumn{2}{c}{CrossNER\_Lit} & \multicolumn{2}{c}{FabNER} & \multicolumn{2}{c}{Mit-Movie} & \multicolumn{2}{c}{MultiNERD} & \multicolumn{2}{c}{TweetNER} & \multicolumn{2}{c}{WikiANN} & \multicolumn{2}{c}{Avg} \\
\cmidrule(lr){2-3} \cmidrule(lr){4-5} \cmidrule(lr){6-7} \cmidrule(lr){8-9} \cmidrule(lr){10-11} \cmidrule(lr){12-13} \cmidrule(lr){14-15} \cmidrule(lr){16-17} \cmidrule(lr){18-19}
& Prec & F1 & Prec & F1 & Prec & F1 & Prec & F1 & Prec & F1 & Prec & F1 & Prec & F1 & Prec & F1 & Prec & F1 \\
\midrule
TIES & 8.33 & 6.26 & 4.56 & 3.82 & 4.69 & 3.82 & 7.00 & 5.55 & 6.02 & 5.14 & 20.93 & 22.87 & 10.73 & 9.94 & 40.27 & 38.99 & 12.82 & 12.05 \\
DARE & 5.62 & 4.48 & 1.89 & 1.51 & 1.41 & 1.28 & 3.54 & 2.93 & 2.62 & 2.17 & 11.63 & 12.55 & 5.61 & 5.23 & 31.44 & 30.52 & 7.97 & 7.58 \\
DELLA & 5.65 & 4.47 & 1.52 & 1.12 & 1.46 & 1.25 & 3.88 & 3.11 & 2.52 & 2.17 & 12.13 & 13.18 & 5.83 & 5.32 & 31.11 & 30.13 & 8.01 & 7.59 \\
Breadcrumbs & 17.04 & 18.15 & 21.56 & 14.90 & 30.82 & 27.09 & 11.00 & 6.06 & 20.80 & 15.73 & 34.03 & 47.72 & 27.07 & 22.90 & 43.06 & 39.36 & 25.67 & 23.99 \\
Task Arithmetic & 16.42 & 17.67 & 22.61 & 14.15 & 30.89 & 25.48 & 11.03 & 5.47 & 20.29 & 14.47 & 36.10 & 49.15 & 26.88 & 22.28 & 44.35 & 40.14 & 26.07 & 23.60 \\
EMR-Merging & 16.34 & 17.90 & 20.97 & 13.02 & 29.12 & 24.81 & 11.70 & 5.80 & 20.47 & 14.53 & 35.65 & 49.08 & 26.26 & 21.25 & 43.28 & 39.11 & 25.47 & 23.19 \\
LoRAHub & 23.85 & 20.59 & 32.76 & 32.44 & 35.56 & 32.72 & 26.45 & 22.48 & 35.62 & 33.12 & 35.46 & 42.15 & 25.12 & 27.27 & 52.70 & 54.12 & 33.44 & 33.11 \\
\midrule
Evo-Merging+ & \textbf{28.62} & \textbf{22.58} & \textbf{42.03} & \textbf{41.03} & \textbf{40.31} & \textbf{37.82} & \textbf{35.71} & \textbf{31.27} & \textbf{56.77} & \textbf{52.34} & \textbf{44.68} & \textbf{51.00} & \textbf{37.41} & \textbf{38.16} & \textbf{63.20} & \textbf{63.40} & \textbf{43.59} & \textbf{42.20} \\
\bottomrule
\end{tabular}
\vspace{-2mm}
\end{table*}


\section{Details of General Model Merging Baselines}
\label{sec:general_baseline_details}

In this section, we provide the detailed descriptions and formulations of the general-purpose model merging methods used in our comparison.

\paragraph{White-box Methods (Parameter-Access Required).}
\begin{itemize}[leftmargin=*, align=left]
    \item \textbf{Task Arithmetic~\cite{ilharco2022editing}.} 
    A simple yet effective method that linearly combines task vectors:
    $\theta_m = \theta_{\text{pre}} + \lambda \sum_t \tau_t$, 
    where $\lambda$ is a scaling factor controlling the strength of task integration.

    \item \textbf{TIES-Merging~\cite{yadav2023ties}.} 
    Addresses interference among tasks through a three-step procedure: \textbf{T}rimming low-magnitude parameters, \textbf{E}lecting a dominant sign across tasks, and \textbf{S}eparately merging parameters that agree in sign:
    $\tau_m = \text{TIES}(\{\tau_t\})$.

    \item \textbf{DARE (Drop and Rescale)~\cite{yu2024languagemodelssupermario}.} 
    Sparsifies task vectors by randomly \textbf{D}ropping a fraction $p$ of delta parameters and \textbf{RE}scaling the remaining ones to preserve the expected magnitude:
    $\delta' = \frac{\delta \odot (1 - m)}{1 - p}$, where $m \sim \text{Bernoulli}(p)$.

    \item \textbf{DELLA-Merging~\cite{deep2024dellamergingreducinginterferencemodel}.} 
    Introduces MAGPRUNE, a magnitude-aware pruning strategy that preferentially drops small-magnitude deltas, followed by sign election and fusion:
    $p(\text{drop}_i) \propto \frac{1}{|\delta_i|}$.

    \item \textbf{EMR-Merging~\cite{huang2024emr}.} 
    Constructs a unified model $\theta_{\text{unified}}$, then recovers task-specific behavior via lightweight \textbf{M}asks ($M_t$) and \textbf{R}escalers ($\lambda_t$):
    $\theta'_t = \lambda_t \cdot (M_t \odot \theta_{\text{unified}})$.

    \item \textbf{Breadcrumbs~\cite{davari2024modelbreadcrumbsscalingmultitask}.} 
    Generates layer-wise sparse “breadcrumb” masks $m_t$ that filter out both negligible and outlier weights from task vectors, enabling cleaner multi-task composition:
    $\theta_m = \theta_{\text{pre}} + \alpha \sum_t (m_t \odot \tau_t)$.

    \item \textbf{TSV-M~\cite{chen2024tsvm}.}
    Uses task singular vectors to identify and merge informative directions in task-vector space, requiring access to adapter parameters.

    \item \textbf{AdaMerging~\cite{yang2023adamerging}.}
    Learns adaptive merging coefficients for different model components using gradient-based optimization on validation data, and is therefore a white-box/gradient-based comparator rather than a black-box API-only method.

    \item \textbf{Robust Merging~\cite{li2024robustmerging}.}
    Improves robustness of parameter merging by reducing the influence of conflicting or noisy update directions, while still assuming direct access to model deltas.
\end{itemize}

\paragraph{Black-box Compatible Method.}
\begin{itemize}[leftmargin=*, align=left]
    \item \textbf{LoRAHub~\cite{huang2024lorahubefficientcrosstaskgeneralization}.} 
    Achieves cross-task generalization by dynamically composing LoRA modules $\{m_i\}$ via gradient-free optimization to learn combination weights $\{w_i\}$: $m_{\text{composed}} = \sum_i w_i m_i$.
\end{itemize}

\section{Algorithm Flow of Evo-Merging}
\label{alg: Algorithm flow of Evo-Merging}
We summarize the procedure of \texttt{Evo-Merging} in Algorithm \ref{alg:Evo-Merging}.
\begin{algorithm}[htbp]
\footnotesize
\caption{Evo-Merging Procedure}
\label{alg:Evo-Merging}
\begin{algorithmic}[1] 
\REQUIRE Collection of LoRA pairs $\{A_i, B_i\}_{i=1..N}$, validation data $\mathcal{D}_{\text{val}}$, regularization parameters $\lambda_1, \lambda_2$.
\ENSURE Merged LoRA pair $(\mathcal{A}_m, \mathcal{B}_m)$.
\STATE

\STATE $\triangleright$ \textbf{Step 1: Sparsity-Based Denoising}
\STATE Set uniform weights $w_i \gets 1/N$ for $i=1..N$.
\STATE $B_{\text{pm}} \gets \sum_{i=1}^{N} w_i \cdot B_i$.
\STATE Use CMA-ES to find optimal $\bm{\alpha}^*$:

    \STATE \quad $\bm{\alpha}^* \gets \underset{\bm{\alpha} \in [0,1]^N}{\text{argmin}} \Big( \mathcal{L}_{\text{CE}}(\mathcal{M}(A_{\text{pm}}(\bm{\alpha}), B_{\text{pm}}), \mathcal{D}_{\text{val}})$
    \STATE \quad \qquad\qquad\qquad\quad ${} + \lambda_1 ||\bm{\alpha}||_1 \Big)$
    \STATE \quad where $A_{\text{pm}}(\bm{\alpha}) = \sum_{i=1}^{N} w_i \cdot \mathbf{S}_{\alpha_i}(A_i)$.
\STATE

\STATE $\triangleright$ \textbf{Step 2: Sign-Aware Scaling}
\STATE Create sparsified LoRA-A matrices:
\FOR{$i = 1, \dots, N$}
    \STATE $A'_i \gets \mathbf{S}_{\alpha_i^*}(A_i)$
\ENDFOR
\STATE Use CMA-ES to find optimal $\bm{\beta}^*$:
    \STATE \quad $\bm{\beta}^* \gets \underset{\bm{\beta} \in \mathbb{R}^N}{\text{argmin}} \Big( \mathcal{L}_{\text{CE}}(\mathcal{M}(\bm{\beta}), \mathcal{D}_{\text{val}})$
    \STATE \quad \qquad\qquad\qquad ${} + \lambda_2 ||\bm{\beta}||_1 \Big)$
    \STATE \quad where $\mathcal{M}(\bm{\beta})$ uses $\mathcal{A}_m = \sum \beta_i A'_i, \mathcal{B}_m = \sum \beta_i B_i$.
\STATE

\STATE $\triangleright$ \textbf{Step 3: Construct Final Merged Model}
\STATE $\mathcal{A}_m \gets \sum_{i=1}^{N} \beta_i^* \cdot A'_i$
\STATE $\mathcal{B}_m \gets \sum_{i=1}^{N} \beta_i^* \cdot B_i$
\end{algorithmic}
\end{algorithm}


\section{Hyperparameter Grid Search for Baselines}
\label{sec:appendix_hyperparams}

This section provides a summary of the hyperparameter grid search conducted for the baseline methods across all datasets. 
This table \ref{tab:grid_search_summary} presents the optimal configurations and their corresponding best performance (Precision and F1-score in \%). The grid search for $\lambda$ was conducted over the set \{0.1, 0.3, 0.5, 0.7, 1.0\}, and for the density $p$ over \{0.5, 0.7, 0.9, 1.0\}. 




\section{Full Per-Task Result Tables}
\label{sec:full_tables}

\subsection{Detailed Per-Task Results (Llama 3.1)}
\paragraph{OOD Full Results (Llama 3.1)}

\begin{table*}[htbp]
\centering
\caption{Full per-task results corresponding to the OOD averages in Table~\ref{tab:method_performance_merged}. \texttt{Evo-Merging} and LoRAHub operate on the full repository; other baselines follow the pre-selection strategy described in Appendix~\ref{sec: Pre-selection Strategy for Baselines in OOD Setting}.}
\label{tab:method_performance_ood_full}
\vspace{-2mm}
\setlength{\tabcolsep}{4.0pt}
\begin{tabular}{l cc cc cc cc cc cc c c}
\toprule
\multirow{3}{*}{{Method}}  & \multicolumn{4}{c}{{NER}} & \multicolumn{4}{c}{{RE}} & \multicolumn{4}{c}{{ET}} & \multicolumn{2}{c}{\multirow{2}{*}{{Avg}}} \\ 
\cmidrule(lr){2-5} \cmidrule(lr){6-9} \cmidrule(lr){10-13}
 & \multicolumn{2}{c}{{Mit-Restaurant}} & \multicolumn{2}{c}{{FindVehicle}} & \multicolumn{2}{c}{{New York Time}} & \multicolumn{2}{c}{{Conll04}} & \multicolumn{2}{c}{{Mit-Movie}} & \multicolumn{2}{c}{{FabNER}} & & \\
\cmidrule(lr){2-3} \cmidrule(lr){4-5} \cmidrule(lr){6-7} \cmidrule(lr){8-9} \cmidrule(lr){10-11} \cmidrule(lr){12-13} \cmidrule(lr){14-15}
  & Prec & F1 & Prec & F1 & Prec & F1 & Prec & F1 & Prec & F1 & Prec & F1 & Prec & F1 \\
\midrule
TIES            & 21.40 & 24.22 & 29.87 & 29.45 & 12.55 & 7.57  & 20.78 & 19.05 & 48.67 & 51.12 & 15.91 & 11.42 & 24.86 & 23.81 \\
DARE (TIES+)    & 10.88 & 14.31 & 24.13 & 23.59 & 3.69  & 2.53  & 10.85 & 8.08  & 44.57 & 48.60 & 15.44 & 10.71 & 18.26 & 17.97 \\
Task Arithmetic & 12.02 & 11.57 & 14.29 & 14.96 & 15.24 & 16.04 & 11.44 & 12.03 & 36.70 & 41.36 & 9.90  & 9.68  & 16.60 & 17.61 \\
Breadcrumbs      & 38.99 & 34.06 & 47.41 & 38.42 & 34.05 & 28.03 & 26.30 & 24.56 & 49.73 & 51.60 & 20.53 & 16.84 & 36.17 & 32.25 \\
Della           & 29.69 & 26.06 & 16.66 & 17.04 & 3.94  & 3.25  & 19.91 & 17.97 & 43.85 & 47.68 & 14.41 & 12.51 & 21.41 & 20.75 \\
EMR-Merging     & 31.89 & 27.46 & 46.92 & 42.25 & 32.49 & 28.94 & 16.45 & 17.10 & 46.57 & 47.35 & 19.26 & 17.86 & 32.26 & 30.16 \\
LoRAHub         & 65.05 & 63.08 & 31.81 & 30.75 & 54.23 & 48.61 & 21.30 & 21.77 & 41.99 & 41.70 & 41.94 & 39.74 & 42.72 & 40.94 \\
Best Adapter    & 37.71 & 35.99 & 39.11 & 38.49 & 45.55 & 38.77 & 11.47 & 11.33 & 46.76 & 46.16 & \textbf{48.45} & \textbf{47.16} & 38.18 & 36.32 \\
LoRA            & 28.18 & 29.65 & \textbf{52.19} & \textbf{52.34} & 52.05 & 44.66 & \textbf{45.42} & 41.60 & 25.20 & 25.38 & 18.74 & 20.01 & 36.96 & 35.61 \\
(IA)$^3$        & 36.18 & 41.64 & 24.98 & 26.05 & 30.54 & 30.70 & 17.01 & 19.97 & 38.80 & 42.17 & 15.38 & 16.70 & 27.15 & 29.54 \\
TSV-M           & 27.88 & 26.04 & 7.71 & 3.38 & 55.26 & 47.28 & 16.67 & 15.39 & 76.75 & 76.65 & 0.61 & 0.53 & 30.81 & 28.21 \\
AdaMerging      & 12.70 & 10.84 & 42.43 & 30.52 & 48.17 & 40.64 & 9.39 & 9.62 & \textbf{78.10} & \textbf{78.75} & 24.45 & 23.94 & 35.87 & 32.39 \\
Robust Merging  & 27.63 & 27.04 & 0.08 & 0.05 & 38.39 & 33.39 & 17.17 & 17.55 & 66.72 & 66.37 & 1.17 & 1.07 & 25.19 & 24.25 \\
\midrule
\textbf{\texttt{Evo-Merging}} & \textbf{66.29} & \textbf{65.99} & 49.61 & 47.64 & \textbf{57.81} & \textbf{51.15} & 41.87 & \textbf{43.00} & 64.11 & 62.55 & 43.10 & 42.47 & \textbf{53.80} & \textbf{52.13} \\
\bottomrule
\end{tabular}
\vspace{-1mm}
\end{table*}


\paragraph{ID Full Results (Llama 3.1)}

\begin{table*}[htbp]
\centering
\caption{Full per-task results corresponding to the ID averages in Table~\ref{tab:method_performance_merged}.}
\label{tab:method_performance_indomain_full}
\vspace{-2mm}
\setlength{\tabcolsep}{1.3pt}
\begin{tabular}{l *{9}{cc}}
\toprule
\multirow{2}{*}{\textbf{Method}}
& \multicolumn{2}{c}{ACE\_2005} & \multicolumn{2}{c}{CrossNER\_AI} & \multicolumn{2}{c}{CrossNER\_Lit} & \multicolumn{2}{c}{FabNER} & \multicolumn{2}{c}{Mit-Movie} & \multicolumn{2}{c}{MultiNERD} & \multicolumn{2}{c}{TweetNER} & \multicolumn{2}{c}{WikiANN} & \multicolumn{2}{c}{Avg} \\
\cmidrule(lr){2-3} \cmidrule(lr){4-5} \cmidrule(lr){6-7} \cmidrule(lr){8-9} \cmidrule(lr){10-11} \cmidrule(lr){12-13} \cmidrule(lr){14-15} \cmidrule(lr){16-17} \cmidrule(lr){18-19}
& Prec & F1 & Prec & F1 & Prec & F1 & Prec & F1 & Prec & F1 & Prec & F1 & Prec & F1 & Prec & F1 & Prec & F1 \\
\midrule
Individual model & 70.45 & 69.37 & 9.51 & 21.23 & 20.05 & 11.83 & 48.61 & 47.09 & 74.27 & 74.38 & 86.15 & 85.55 & 62.24 & 58.26 & 81.33 & 80.82 & 56.58 & 56.07 \\
\hline
TIES & 16.16 & 13.11 & 18.99 & 18.93 & 27.27 & 26.47 & 21.02 & 18.91 & \textbf{54.77} & 51.24 & 40.70 & 48.99 & 34.05 & 33.96 & 55.30 & 56.47 & 33.53 & 33.51 \\
DARE & 12.72 & 13.50 & 20.40 & 18.42 & 25.97 & 25.13 & 19.11 & 18.47 & 53.12 & 51.67 & 51.58 & 47.14 & 32.97 & 32.91 & 54.13 & 55.58 & 33.75 & 32.85 \\
DELLA & 15.73 & 12.74 & 17.72 & 18.02 & 25.52 & 25.16 & 19.98 & 18.55 & 54.65 & \textbf{51.94} & 38.69 & 47.11 & 31.99 & 32.30 & 54.61 & 56.15 & 32.36 & 32.75 \\
Breadcrumbs & 19.74 & 15.34 & 21.33 & 20.03 & 31.87 & 29.00 & 21.18 & 17.43 & 51.46 & 46.91 & 48.44 & 55.53 & 37.26 & 33.65 & 56.56 & 56.82 & 35.98 & 34.34 \\
Task Arithmetic & 17.37 & 13.40 & 20.12 & 18.90 & 30.55 & 28.38 & 20.86 & 17.63 & 53.59 & 49.27 & 45.43 & 53.04 & 35.60 & 33.47 & 56.33 & 56.93 & 34.98 & 33.88 \\
LoRAHub & 28.36 & 25.18 & 33.67 & 33.69 & 33.23 & 31.51 & 26.04 & 22.32 & 38.90 & 36.48 & 36.61 & 43.89 & 26.99 & 30.04 & 53.78 & 55.32 & 34.70 & 34.80 \\
EMR-Merging & 19.08 & 15.10 & 16.99 & 16.16 & 26.09 & 23.12 & 23.70 & 19.76 & 52.09 & 46.67 & 46.65 & 54.16 & 37.91 & 34.76 & 59.01 & 59.49 & 35.19 & 33.65 \\
TSV-M & \textbf{43.03} & \textbf{37.32} & 27.57 & 25.41 & 21.40 & 18.90 & \textbf{31.03} & \textbf{27.29} & 52.85 & 49.76 & 30.19 & 36.37 & 31.78 & 33.47 & 55.09 & 55.54 & 36.62 & 35.51 \\
AdaMerging & 13.37 & 8.19 & \textbf{38.30} & 29.00 & 27.80 & 23.24 & 24.77 & 12.88 & 34.90 & 26.54 & \textbf{74.53} & \textbf{74.88} & 31.60 & 26.73 & 59.44 & 57.97 & 38.09 & 32.43 \\
Robust Merging & 20.72 & 15.73 & 37.60 & \textbf{34.93} & 32.80 & 30.18 & 24.78 & 18.92 & 46.28 & 40.51 & 45.05 & 52.06 & \textbf{40.01} & \textbf{40.10} & \textbf{61.76} & \textbf{61.57} & 38.63 & 36.75 \\
\midrule
\textbf{\texttt{Evo-Merging}} & 26.31 & 23.46 & 34.18 & 34.61 & \textbf{33.95} & \textbf{32.32} & 27.15 & 23.27 & 43.58 & 39.94 & 52.82 & 56.59 & 36.34 & 35.08 & 58.40 & 58.98 & \textbf{39.09} & \textbf{38.03} \\
\bottomrule
\end{tabular}
\vspace{-1mm}
\end{table*}


\subsection{Ablation Full Results}

\begin{table*}[htbp]
\centering
\caption{Full per-task ablation study corresponding to the average ablation results in Table~\ref{tab:detailed_ablation_final_v2}. [$\downarrow$] means the performance drop relative to our full model.}
\label{tab:detailed_ablation_final_v2_full}
\vspace{-2mm}
\setlength{\tabcolsep}{0.5mm}
\begin{tabular}{lcccccccccccc cc} 
\toprule
\multirow{2}{*}{Method}  & \multicolumn{2}{c}{NER\_Mit-Rest} & \multicolumn{2}{c}{NER\_FindVeh} & \multicolumn{2}{c}{RE\_NYT} & \multicolumn{2}{c}{RE\_Conll04} & \multicolumn{2}{c}{ET\_Mit-Movie} & \multicolumn{2}{c}{ET\_FabNER} & \multicolumn{2}{c}{Avg} \\
\cmidrule(lr){2-3} \cmidrule(lr){4-5} \cmidrule(lr){6-7} \cmidrule(lr){8-9} \cmidrule(lr){10-11} \cmidrule(lr){12-13} \cmidrule(lr){14-15}
 & Prec & F1 & Prec & F1 & Prec & F1 & Prec & F1 & Prec & F1 & Prec & F1 & Prec & F1 \\
\midrule
\textbf{\texttt{Evo-Merging}} & \textbf{66.29} & \textbf{65.99} & \textbf{49.61} & \textbf{47.64} & \textbf{57.81} & \textbf{51.15} & \textbf{41.87} & \textbf{43.00} & \textbf{64.11} & \textbf{62.55} & \textbf{43.10} & \textbf{42.47} & \textbf{53.80} & \textbf{52.13} \\
\midrule
w/o Denoising \& Scaling & 28.01 & 33.10 & 44.14 & 32.68 & 2.35 & 1.94 & 7.36 & 5.37 & 43.70 & 39.67 & 14.10 & 11.75 & 23.28 [$\downarrow$ 30.5] & 20.75 [$\downarrow$ 31.4] \\
w/o Denoising & 65.05 & 63.08 & 31.81 & 30.75 & 54.23 & 48.61 & 21.30 & 21.77 & 41.99 & 41.70 & 41.94 & 39.74 & 42.72 [$\downarrow$ 11.1] & 40.94 [$\downarrow$ 11.2] \\
w/o Scaling & 39.90 & 37.41 & 48.78 & 41.35 & 2.27 & 2.23 & 7.14 & 6.38 & 46.10 & 46.52 & 14.10 & 11.75 & 26.38 [$\downarrow$ 27.4] & 24.27 [$\downarrow$ 27.9] \\
w/o Sign-flipped & 42.04 & 37.70 & 46.61 & 41.07 & 45.45 & 38.38 & 35.21 & 32.91 & 62.19 & 60.47 & 34.77 & 29.80 & 44.38 [$\downarrow$ 9.4] & 40.06 [$\downarrow$ 12.1] \\
\bottomrule
\end{tabular}
\vspace{-1mm}
\end{table*}
